\section{Линейная зависимость и независимость системы векторов. Критерий линейной зависимости. Свойства ЛЗ и ЛНЗ векторов. Теорема о добавлении вектора}

\begin{definition}
Система векторов $a_1, a_2, \dots, a_k$ из линейного пространства $V$ называется \textbf{линейно зависимой (ЛЗ)}, если существуют числа $\lambda_1, \lambda_2, \dots, \lambda_k \in F$, не все равные нулю одновременно, такие что:
\[
\lambda_1 a_1 + \lambda_2 a_2 + \dots + \lambda_k a_k = 0.
\]
Если же равенство $\lambda_1 a_1 + \dots + \lambda_k a_k = 0$ выполняется только при $\lambda_1 = \lambda_2 = \dots = \lambda_k = 0$, то система векторов называется \textbf{линейно независимой (ЛНЗ)}.
\end{definition}

\begin{theorem}[Критерий линейной зависимости]
Система векторов $a_1, \dots, a_k$ ($k \ge 2$) линейно зависима тогда и только тогда, когда хотя бы один из этих векторов линейно выражается через остальные.
\end{theorem}

\begin{proof}
\begin{itemize}
    \item \textbf{Необходимость ($\Rightarrow$):} Пусть система ЛЗ. Тогда $\exists \lambda_1, \dots, \lambda_k$, не все равные нулю, такие что $\sum_{i=1}^k \lambda_i a_i = 0$. Пусть, для определенности, $\lambda_k \neq 0$. Тогда:
    \[
    \lambda_k a_k = -\lambda_1 a_1 - \dots - \lambda_{k-1} a_{k-1} \quad \Rightarrow \quad a_k = \left(-\frac{\lambda_1}{\lambda_k}\right) a_1 + \dots + \left(-\frac{\lambda_{k-1}}{\lambda_k}\right) a_{k-1},
    \]
    то есть $a_k$ линейно выражается через остальные векторы.
    \item \textbf{Достаточность ($\Leftarrow$):} Пусть, например, $a_k$ линейно выражается через остальные: $a_k = \mu_1 a_1 + \dots + \mu_{k-1} a_{k-1}$. Перенесем $a_k$ в левую часть:
    \[
    \mu_1 a_1 + \dots + \mu_{k-1} a_{k-1} + (-1) a_k = 0.
    \]
    Коэффициент при $a_k$ равен $-1 \neq 0$, следовательно, найдена нетривиальная линейная комбинация, равная нулю. Система ЛЗ.
\end{itemize}
\hfill$\blacksquare$
\end{proof}

\begin{theorem}[Свойства линейно зависимых и независимых систем]
\begin{enumerate}
    \item Если система векторов содержит нулевой вектор, то она линейно зависима.
    \item Если система векторов содержит линейно зависимую подсистему, то вся система линейно зависима.
    \item Если система векторов линейно независима, то любая ее подсистема также линейно независима.
\end{enumerate}
\end{theorem}

\begin{proof}
\begin{enumerate}
    \item Пусть $a_1 = 0$. Тогда $1 \cdot a_1 + 0 \cdot a_2 + \dots + 0 \cdot a_k = 0$. Коэффициент при $a_1$ не равен нулю, значит, система ЛЗ.
    \item Пусть подсистема $a_1, \dots, a_m$ ($m < k$) ЛЗ. Тогда $\exists \lambda_1, \dots, \lambda_m$, не все нули, такие что $\sum_{i=1}^m \lambda_i a_i = 0$. Дополним эту комбинацию нулевыми коэффициентами для остальных векторов: $\sum_{i=1}^m \lambda_i a_i + 0 \cdot a_{m+1} + \dots + 0 \cdot a_k = 0$. Эта комбинация нетривиальна, значит, вся система ЛЗ.
    \item Доказывается от противного, как контрапозиция свойства 2. Если бы подсистема была ЛЗ, то по свойству 2 вся система была бы ЛЗ, что противоречит условию.
\end{enumerate}
\hfill$\blacksquare$
\end{proof}

\begin{theorem}[О добавлении вектора]
Пусть система векторов $a_1, \dots, a_k$ линейно независима, а система $a_1, \dots, a_k, b$ линейно зависима. Тогда вектор $b$ линейно выражается через $a_1, \dots, a_k$.
\end{theorem}

\begin{proof}
Так как система $a_1, \dots, a_k, b$ линейно зависима, существуют числа $\lambda_1, \dots, \lambda_k, \mu$, не все равные нулю, такие что:
\[
\lambda_1 a_1 + \dots + \lambda_k a_k + \mu b = 0.
\]
Предположим, что $\mu = 0$. Тогда $\lambda_1 a_1 + \dots + \lambda_k a_k = 0$. Но система $a_1, \dots, a_k$ по условию линейно независима, следовательно, все $\lambda_i = 0$. Это означает, что все коэффициенты в комбинации равны нулю, что противоречит линейной зависимости системы с вектором $b$. 
Значит, наше предположение неверно, и $\mu \neq 0$. Тогда мы можем выразить $b$:
\[
b = \left(-\frac{\lambda_1}{\mu}\right) a_1 + \dots + \left(-\frac{\lambda_k}{\mu}\right) a_k.
\]
\hfill$\blacksquare$
\end{proof}
